% !TeX root = part4.tex

\documentclass[../final.tex]{subfiles}

\begin{document}


% ============================================================
% Семинар 4
% ============================================================

\practice{4}{26.09.2026}{Колебания решётки. Продолжение}


% ============================================================
% Продолжение задачи о дефектной пружине
% ============================================================

\rtheory{Продолжение задачи о дефектной пружине}

Продолжим исследование локализованного колебательного состояния
в двухатомной цепочке с чередующимися массами $m$ и $M$.

Вдали от дефекта все пружины имеют жёсткость $\alpha$,
а одна из связей имеет жёсткость $\alpha_0$.

Для локализованного состояния амплитуда колебаний
экспоненциально убывает при удалении от дефекта,
поэтому вводим декремент затухания
%
\begin{equation*}
    \chi>0.
\end{equation*}


% ------------------------------------------------------------
% Уравнение для локализованной моды
% ------------------------------------------------------------

\subsection*{Уравнение для локализованной моды}

Из предыдущего семинара для области вдали от дефекта
было получено уравнение
%
\begin{equation*}
    \omega^4
    -
    \frac{2\alpha(m+M)}{mM}\omega^2
    +
    \frac{4\alpha^2}{mM}
    \cosh^2(\chi a)
    =
    0.
\end{equation*}
%

Будем рассматривать состояние,
локализованное слабо, то есть
%
\begin{equation*}
    \chi a\ll1.
\end{equation*}


% ------------------------------------------------------------
% Разложение cosh
% ------------------------------------------------------------

\subsection*{Разложение при малом $\chi a$}

Напомним,
%
\begin{equation*}
    \cosh x
    =
    \frac{e^x+e^{-x}}{2}.
\end{equation*}
%

При малом $x$
%
\begin{equation*}
    e^x
    =
    1+x+\frac{x^2}{2}+\ldots,
\end{equation*}
%
\begin{equation*}
    e^{-x}
    =
    1-x+\frac{x^2}{2}+\ldots.
\end{equation*}
%

Поэтому
%
\begin{equation*}
    \cosh x
    \simeq
    1+\frac{x^2}{2}.
\end{equation*}
%

Возводя в квадрат,
%
\begin{align*}
    \cosh^2x
    &\simeq
    \left(
        1+\frac{x^2}{2}
    \right)^2
    \\
    &=
    1+x^2+\frac{x^4}{4}.
\end{align*}
%

В рассматриваемом приближении член порядка $x^4$
отбрасываем, поэтому
%
\begin{equation*}
    \cosh^2x
    \simeq
    1+x^2.
\end{equation*}
%

Следовательно,
%
\begin{equation*}
    \cosh^2(\chi a)
    \simeq
    1+(\chi a)^2.
\end{equation*}
%

Уравнение для частоты принимает вид
%
\begin{equation*}
    \omega^4
    -
    \frac{2\alpha(m+M)}{mM}\omega^2
    +
    \frac{4\alpha^2}{mM}
    \left[
        1+(\chi a)^2
    \right]
    =
    0.
\end{equation*}


% ------------------------------------------------------------
% Дискриминант
% ------------------------------------------------------------

\subsection*{Дискриминант}

Рассматриваем это уравнение как квадратное
относительно
%
\begin{equation*}
    x=\omega^2.
\end{equation*}
%

Тогда
%
\begin{equation*}
    x^2
    -
    \frac{2\alpha(m+M)}{mM}x
    +
    \frac{4\alpha^2}{mM}
    \left[
        1+(\chi a)^2
    \right]
    =
    0.
\end{equation*}
%

Дискриминант равен
%
\begin{equation*}
    D
    =
    \left[
        \frac{2\alpha(m+M)}{mM}
    \right]^2
    -
    \frac{16\alpha^2}{mM}
    \left[
        1+(\chi a)^2
    \right].
\end{equation*}
%

Удобно рассмотреть величину $D/4$:
%
\begin{equation*}
    \frac{D}{4}
    =
    \alpha^2
    \frac{(m+M)^2}{(mM)^2}
    -
    \frac{4\alpha^2}{mM}
    \left[
        1+(\chi a)^2
    \right].
\end{equation*}
%

Разделим постоянную часть и поправку:
%
\begin{equation*}
    \frac{D}{4}
    =
    \alpha^2
    \left[
        \frac{(m+M)^2}{(mM)^2}
        -
        \frac{4}{mM}
    \right]
    -
    \frac{4\alpha^2}{mM}
    (\chi a)^2.
\end{equation*}
%

Приводим первые два слагаемых к общему знаменателю:
%
\begin{align*}
    \frac{(m+M)^2}{(mM)^2}
    -
    \frac{4}{mM}
    &=
    \frac{
        (m+M)^2-4mM
    }{
        (mM)^2
    }
    \\
    &=
    \frac{
        m^2-2mM+M^2
    }{
        (mM)^2
    }
    \\
    &=
    \frac{
        (M-m)^2
    }{
        (mM)^2
    }.
\end{align*}
%

Следовательно,
%
\begin{equation*}
    \frac{D}{4}
    =
    \alpha^2
    \frac{(M-m)^2}{(mM)^2}
    -
    \frac{4\alpha^2}{mM}
    (\chi a)^2.
\end{equation*}
%

Выносим первый множитель:
%
\begin{equation*}
    \frac{D}{4}
    =
    \alpha^2
    \frac{(M-m)^2}{(mM)^2}
    \left[
        1
        -
        \frac{4mM}{(M-m)^2}
        (\chi a)^2
    \right].
\end{equation*}
%

Поэтому
%
\begin{equation*}
    \frac{\sqrt D}{2}
    =
    \alpha
    \frac{|M-m|}{mM}
    \sqrt{
        1
        -
        \frac{4mM}{(M-m)^2}
        (\chi a)^2
    }.
\end{equation*}


% ------------------------------------------------------------
% Частоты
% ------------------------------------------------------------

\subsection*{Частоты локализованного состояния}

Корни квадратного уравнения имеют вид
%
\begin{equation*}
    \omega_\pm^2
    =
    \alpha
    \frac{m+M}{mM}
    \pm
    \alpha
    \frac{|M-m|}{mM}
    \sqrt{
        1
        -
        \frac{4mM}{(M-m)^2}
        (\chi a)^2
    }.
\end{equation*}
%

Поскольку
%
\begin{equation*}
    \chi a\ll1,
\end{equation*}
%
используем разложение
%
\begin{equation*}
    \sqrt{1-x}
    \simeq
    1-\frac{x}{2}.
\end{equation*}
%

Получаем
%
\begin{equation*}
    \omega_\pm^2
    \simeq
    \alpha
    \frac{m+M}{mM}
    \pm
    \alpha
    \frac{|M-m|}{mM}
    \left[
        1
        -
        \frac{2mM}{(M-m)^2}
        (\chi a)^2
    \right].
\end{equation*}


% ------------------------------------------------------------
% M > m
% ------------------------------------------------------------

\subsection*{Случай $M>m$}

Далее будем считать
%
\begin{equation*}
    M>m.
\end{equation*}
%

Тогда
%
\begin{equation*}
    |M-m|
    =
    M-m.
\end{equation*}
%

Для верхнего корня имеем
%
\begin{align*}
    \omega_+^2
    &\simeq
    \alpha
    \frac{m+M}{mM}
    +
    \alpha
    \frac{M-m}{mM}
    -
    \frac{2\alpha}{M-m}
    (\chi a)^2
    \\
    &=
    \frac{
        \alpha(m+M+M-m)
    }{
        mM
    }
    -
    \frac{2\alpha}{M-m}
    (\chi a)^2
    \\
    &=
    \frac{2\alpha}{m}
    -
    \frac{2\alpha}{M-m}
    (\chi a)^2.
\end{align*}
%

Таким образом,
%
\begin{equation*}
    \omega_+^2
    \simeq
    \frac{2\alpha}{m}
    -
    \frac{2\alpha}{M-m}
    (\chi a)^2.
\end{equation*}
%

Для нижнего корня
%
\begin{align*}
    \omega_-^2
    &\simeq
    \alpha
    \frac{m+M}{mM}
    -
    \alpha
    \frac{M-m}{mM}
    +
    \frac{2\alpha}{M-m}
    (\chi a)^2
    \\
    &=
    \frac{
        \alpha(m+M-M+m)
    }{
        mM
    }
    +
    \frac{2\alpha}{M-m}
    (\chi a)^2
    \\
    &=
    \frac{2\alpha}{M}
    +
    \frac{2\alpha}{M-m}
    (\chi a)^2.
\end{align*}
%

Следовательно,
%
\begin{equation*}
    \omega_-^2
    \simeq
    \frac{2\alpha}{M}
    +
    \frac{2\alpha}{M-m}
    (\chi a)^2.
\end{equation*}


% ------------------------------------------------------------
% Щель
% ------------------------------------------------------------

\subsection*{Смещение частот в запрещённую зону}

Обозначим края запрещённой зоны идеальной цепочки:
%
\begin{equation*}
    \omega_1^2
    =
    \frac{2\alpha}{M},
    \qquad
    \omega_2^2
    =
    \frac{2\alpha}{m}.
\end{equation*}
%

Тогда полученные частоты можно записать как
%
\begin{equation*}
    \omega_-^2
    =
    \omega_1^2
    +
    \Delta\omega^2,
\end{equation*}
%
\begin{equation*}
    \omega_+^2
    =
    \omega_2^2
    -
    \Delta\omega^2,
\end{equation*}
%
где
%
\begin{equation*}
    \Delta\omega^2
    =
    \frac{2\alpha}{M-m}
    (\chi a)^2.
\end{equation*}
%

Таким образом, обе частоты смещаются
от краёв разрешённых ветвей
\rconcept{внутрь запрещённой зоны}.


% ------------------------------------------------------------
% РИСУНОК
% ------------------------------------------------------------

\begin{rbox}[left=18pt,right=18pt,top=14pt,bottom=12pt]{[]}
    \centering

    \begin{tikzpicture}[
        x=1.1cm,
        y=1.0cm,
        >=Latex,
        every node/.style={text=rtext}
    ]

        % Оси
        \draw[
            ->,
            draw=rtext,
            line width=0.9pt
        ]
        (-4.2,0)
        --
        (4.3,0)
        node[right] {$k$};

        \draw[
            ->,
            draw=rtext,
            line width=0.9pt
        ]
        (0,-0.15)
        --
        (0,5.0)
        node[above] {$\omega$};

        % Границы зоны Бриллюэна
        \draw[
            dashed,
            draw=rtext,
            opacity=0.55
        ]
        (-3.5,0)
        --
        (-3.5,4.5);

        \draw[
            dashed,
            draw=rtext,
            opacity=0.55
        ]
        (3.5,0)
        --
        (3.5,4.5);

        \node[below=5pt] at (-3.5,0)
            {$-\dfrac{\pi}{2a}$};

        \node[below=5pt] at (3.5,0)
            {$\dfrac{\pi}{2a}$};

        % Акустическая ветвь
        \draw[
            draw=rc3,
            line width=1.3pt
        ]
        (-3.5,2.0)
        .. controls (-2.2,1.0) and (-0.8,0.15) ..
        (0,0)
        .. controls (0.8,0.15) and (2.2,1.0) ..
        (3.5,2.0);

        % Оптическая ветвь
        \draw[
            draw=p3,
            line width=1.3pt
        ]
        (-3.5,3.3)
        .. controls (-2.1,3.75) and (-0.8,4.15) ..
        (0,4.25)
        .. controls (0.8,4.15) and (2.1,3.75) ..
        (3.5,3.3);

        % Уровни локализованных состояний
        \draw[
            draw=red3,
            line width=1.1pt
        ]
        (-1.0,2.35)
        --
        (1.0,2.35);

        \draw[
            draw=red3,
            line width=1.1pt
        ]
        (-1.0,3.00)
        --
        (1.0,3.00);

        \node[
            right,
            text=red3
        ]
        at (1.05,2.35)
        {$\omega_-^2$};

        \node[
            right,
            text=red3
        ]
        at (1.05,3.00)
        {$\omega_+^2$};

        % Границы щели
        \draw[
            dashed,
            draw=rc3,
            opacity=0.65
        ]
        (0,2.0)
        --
        (3.9,2.0);

        \draw[
            dashed,
            draw=p3,
            opacity=0.65
        ]
        (0,3.3)
        --
        (3.9,3.3);

        \node[
            right,
            text=rc3
        ]
        at (3.55,2.0)
        {$\omega_1^2$};

        \node[
            right,
            text=p3
        ]
        at (3.55,3.3)
        {$\omega_2^2$};

    \end{tikzpicture}

    \captionsetup{
        font=footnotesize,
        labelfont=bf,
        skip=5pt
    }

    \captionof{figure}{
        Локализованные состояния в запрещённой зоне.
        Частота $\omega_-^2$ отщепляется вверх
        от акустической ветви,
        а $\omega_+^2$ --- вниз от оптической.
    }

    \label{fig:defect_gap_states}
\end{rbox}


% ============================================================
% Слабый дефект
% ============================================================

\subsection*{Слабое изменение жёсткости}

Вернёмся к дефектной пружине.

Для центральной области имеем условие
%
\begin{equation*}
    mM\omega^4
    -
    (\alpha_0+\alpha)(m+M)\omega^2
    +
    (\alpha_0+\alpha)^2
    -
    \left(
        \alpha_0-\alpha e^{-2\chi a}
    \right)^2
    =
    0.
\end{equation*}
%

Вдали от дефекта выполняется
%
\begin{equation*}
    mM\omega^4
    -
    2\alpha(m+M)\omega^2
    +
    2\alpha^2
    \left[
        1+\cosh(2\chi a)
    \right]
    =
    0.
\end{equation*}
%

Рассмотрим слабый дефект:
%
\begin{equation*}
    \alpha_0
    =
    \alpha+\gamma,
    \qquad
    |\gamma|\ll\alpha.
\end{equation*}
%

Здесь
%
\begin{equation*}
    \gamma
    =
    \alpha_0-\alpha
\end{equation*}
%
--- малое изменение жёсткости дефектной пружины.


% ------------------------------------------------------------
% Вычитание уравнений
% ------------------------------------------------------------

\subsection*{Вычитание двух уравнений}

Подставим
%
\begin{equation*}
    \alpha_0
    =
    \alpha+\gamma
\end{equation*}
%
в уравнение центральной ячейки.

Тогда
%
\begin{equation*}
    \alpha_0+\alpha
    =
    2\alpha+\gamma,
\end{equation*}
%
а
%
\begin{equation*}
    \alpha_0-\alpha e^{-2\chi a}
    =
    \alpha+\gamma-\alpha e^{-2\chi a}.
\end{equation*}
%

Вычтем уравнение для идеальной области
из уравнения около дефекта.

Члены $mM\omega^4$ сокращаются, и остаётся
%
\begin{align*}
    &
    -\gamma(m+M)\omega^2
    +
    (2\alpha+\gamma)^2
    \\
    &\qquad
    -
    \left(
        \alpha+\gamma-\alpha e^{-2\chi a}
    \right)^2
    -
    2\alpha^2
    \left[
        1+\cosh(2\chi a)
    \right]
    =
    0.
\end{align*}


% ------------------------------------------------------------
% Разложения
% ------------------------------------------------------------

\subsection*{Разложение по малым параметрам}

Используем
%
\begin{equation*}
    e^{-2\chi a}
    \simeq
    1
    -
    2\chi a
    +
    2(\chi a)^2,
\end{equation*}
%
а также
%
\begin{equation*}
    \cosh(2\chi a)
    \simeq
    1
    +
    2(\chi a)^2.
\end{equation*}
%

Кроме того,
%
\begin{equation*}
    (2\alpha+\gamma)^2
    =
    4\alpha^2
    +
    4\alpha\gamma
    +
    \gamma^2.
\end{equation*}
%

Рассмотрим выражение
%
\begin{equation*}
    \alpha+\gamma-\alpha e^{-2\chi a}.
\end{equation*}
%

После подстановки разложения экспоненты
%
\begin{align*}
    \alpha+\gamma-\alpha e^{-2\chi a}
    &\simeq
    \alpha+\gamma
    -
    \alpha
    \left[
        1-2\chi a+2(\chi a)^2
    \right]
    \\
    &=
    \gamma
    +
    2\alpha\chi a
    -
    2\alpha(\chi a)^2.
\end{align*}
%

В низшем порядке по малым параметрам
%
\begin{equation*}
    \left(
        \alpha+\gamma-\alpha e^{-2\chi a}
    \right)^2
    \simeq
    4\alpha^2(\chi a)^2.
\end{equation*}
%

Также
%
\begin{align*}
    2\alpha^2
    \left[
        1+\cosh(2\chi a)
    \right]
    &\simeq
    2\alpha^2
    \left[
        2+2(\chi a)^2
    \right]
    \\
    &=
    4\alpha^2
    +
    4\alpha^2(\chi a)^2.
\end{align*}
%

Отбрасывая члены более высокого порядка,
в частности $\gamma^2$,
получаем
%
\begin{align*}
    &
    -\gamma(m+M)\omega^2
    +
    4\alpha^2
    +
    4\alpha\gamma
    \\
    &\qquad
    -
    4\alpha^2(\chi a)^2
    -
    4\alpha^2
    -
    4\alpha^2(\chi a)^2
    =
    0.
\end{align*}
%

После сокращения
%
\begin{equation*}
    -\gamma(m+M)\omega^2
    +
    4\alpha\gamma
    -
    8\alpha^2(\chi a)^2
    =
    0.
\end{equation*}
%

Следовательно,
%
\begin{equation*}
    8\alpha^2(\chi a)^2
    =
    \gamma
    \left[
        4\alpha
        -
        (m+M)\omega^2
    \right].
\end{equation*}


\remark[Замечание к рукописной записи.]
На листе рядом с этим преобразованием стоит пометка
«Ошибка?».

При последовательном разложении двух последних слагаемых
по малому параметру $\chi a$
коэффициент перед $\alpha^2(\chi a)^2$
получается равным $8$.


% ============================================================
% Верхний край щели
% ============================================================

\subsection*{Состояние около верхнего края запрещённой зоны}

Для состояния,
отщепляющегося от оптической ветви,
можно в правой части использовать
нулевое приближение
%
\begin{equation*}
    \omega^2
    \simeq
    \omega_2^2
    =
    \frac{2\alpha}{m}.
\end{equation*}
%

Тогда
%
\begin{equation*}
    8\alpha^2(\chi a)^2
    =
    \gamma
    \left[
        4\alpha
        -
        (m+M)
        \frac{2\alpha}{m}
    \right].
\end{equation*}
%

Выносим $2\alpha$:
%
\begin{equation*}
    8\alpha^2(\chi a)^2
    =
    2\alpha\gamma
    \left[
        2
        -
        \frac{m+M}{m}
    \right].
\end{equation*}
%

В скобках
%
\begin{align*}
    2-\frac{m+M}{m}
    &=
    \frac{2m-m-M}{m}
    \\
    &=
    \frac{m-M}{m}.
\end{align*}
%

Поэтому
%
\begin{equation*}
    8\alpha^2(\chi a)^2
    =
    2\alpha\gamma
    \frac{m-M}{m}.
\end{equation*}
%

Сокращая на $2\alpha$,
получаем
%
\begin{equation*}
    4\alpha(\chi a)^2
    =
    \gamma
    \frac{m-M}{m}.
\end{equation*}
%

Отсюда
%
\begin{equation*}
    (\chi a)^2
    =
    \frac{\gamma}{4\alpha}
    \frac{m-M}{m}.
\end{equation*}
%

Следовательно,
%
\begin{equation*}
    \chi
    =
    \frac{1}{2a}
    \sqrt{
        \frac{\gamma}{\alpha}
        \frac{m-M}{m}
    }.
\end{equation*}
%

При $M>m$ имеем
%
\begin{equation*}
    m-M<0.
\end{equation*}
%

Для действительного положительного декремента $\chi$
необходимо
%
\begin{equation*}
    \gamma<0.
\end{equation*}
%

То есть
%
\begin{equation*}
    \alpha_0<\alpha.
\end{equation*}
%

Следовательно, ослабление дефектной связи
создаёт локализованное состояние,
которое отщепляется
\rresult{вниз от оптической ветви в запрещённую зону}.


% ============================================================
% Нижний край щели
% ============================================================

\subsection*{Состояние около нижнего края запрещённой зоны}

Для состояния,
отщепляющегося от акустической ветви,
берём
%
\begin{equation*}
    \omega^2
    \simeq
    \omega_1^2
    =
    \frac{2\alpha}{M}.
\end{equation*}
%

Тогда
%
\begin{equation*}
    8\alpha^2(\chi a)^2
    =
    \gamma
    \left[
        4\alpha
        -
        (m+M)
        \frac{2\alpha}{M}
    \right].
\end{equation*}
%

Выносим $2\alpha$:
%
\begin{equation*}
    8\alpha^2(\chi a)^2
    =
    2\alpha\gamma
    \left[
        2
        -
        \frac{m+M}{M}
    \right].
\end{equation*}
%

При этом
%
\begin{align*}
    2-\frac{m+M}{M}
    &=
    \frac{2M-m-M}{M}
    \\
    &=
    \frac{M-m}{M}.
\end{align*}
%

Следовательно,
%
\begin{equation*}
    8\alpha^2(\chi a)^2
    =
    2\alpha\gamma
    \frac{M-m}{M}.
\end{equation*}
%

Отсюда
%
\begin{equation*}
    (\chi a)^2
    =
    \frac{\gamma}{4\alpha}
    \frac{M-m}{M},
\end{equation*}
%
и
%
\begin{equation*}
    \chi
    =
    \frac{1}{2a}
    \sqrt{
        \frac{\gamma}{\alpha}
        \frac{M-m}{M}
    }.
\end{equation*}
%

При $M>m$ действительное значение $\chi$
требует
%
\begin{equation*}
    \gamma>0,
\end{equation*}
%
то есть
%
\begin{equation*}
    \alpha_0>\alpha.
\end{equation*}
%

Следовательно, увеличение жёсткости дефектной пружины
создаёт локализованное состояние,
которое отщепляется
\rresult{вверх от акустической ветви в запрещённую зону}.


% ============================================================
% Итог
% ============================================================

\subsection*{Физический результат}

При
%
\begin{equation*}
    M>m
\end{equation*}
%
запрещённая зона идеальной двухатомной цепочки
ограничена частотами
%
\begin{equation*}
    \omega_1^2
    =
    \frac{2\alpha}{M},
    \qquad
    \omega_2^2
    =
    \frac{2\alpha}{m}.
\end{equation*}
%

Дефект одной связи создаёт локализованное состояние
внутри этой зоны.

Для
%
\begin{equation*}
    \alpha_0<\alpha
\end{equation*}
%
локализованное состояние отщепляется
от оптической ветви вниз.

Для
%
\begin{equation*}
    \alpha_0>\alpha
\end{equation*}
%
локализованное состояние отщепляется
от акустической ветви вверх.

При этом
%
\begin{equation*}
    u_n,\;v_n
    \propto
    e^{-2\chi a|n|},
\end{equation*}
%
поэтому величина $\chi$ характеризует
степень пространственной локализации состояния:
чем больше $\chi$, тем быстрее затухает амплитуда
при удалении от дефекта.


\end{document}